\documentclass{article}

\usepackage{luatexja} % 日本語用（LuaLaTeXならこれ必須）
\title{研究説明}
\author{高久}
\date{\today}

\begin{document}

\maketitle

\section{提案手法}

我々はプロンプトによってLLMの出力を制限する手法を知っているが、プロンプトによる制約には限界がある。\\
LLMは自由に出力できるため、物理制御や状態依存タスクでは誤った値がそのまま実行されるリスクがある。\\

本研究の本質は、LLMの出力をそのまま実行するのではなく、\textbf{制約空間に写してから実行する}ことである。\\  
そのために、出力を一度制約付きの構造に変換し、検証・補正してから実行する仕組みを提案する。\\
例えば、以下のように型と制約を定義する：

\begin{verbatim}
{
  "action": "rotate",
  "angle": "int(0,90)"
}
\end{verbatim}

LLMは候補生成のみを行い、実際の制御はプロトコル側で行う。\\
また、状態は構造的に管理し、実行結果はフィードバックとして次の入力に反映される。

本手法は、状態とフィードバックを持つ状態遷移系として捉えることができる。

\subsection{限界}

LLMの確率的性質により、ハルシネーション自体を完全に防ぐことはできない。  
また、本手法は意味理解を保証するものではなく、実行可能性を制約によって担保するアプローチである。

\section{実験}

IoT環境における障害物回避タスクを用い、制約違反率とタスク成功率を評価する。

\begin{table}[h]
\centering
\begin{tabular}{lcccc}
\hline
指標 & LLM単体 & LLM+RAG & LLM+ECS & 提案手法 \\
\hline
制約違反率 & 高い & 高い & 中 & 低い \\
タスク成功率 & 低い & 低い & 中 & 高い \\
実行時間 & 2m & 2m & 1m & 許容範囲 \\
\hline
\end{tabular}
\caption{各手法の比較}
\end{table}

これらの結果が得られた場合、本手法により安定性向上が示唆される。

\section{位置づけ}

本手法は言語処理そのものではなく、システム工学的観点からAIの制御を扱うアプローチである。

\section{想定される質問への回答}

\subsection{「それRAGでよくない？」}

RAGは情報補完には有効だが、制約違反そのものを防ぐことはできない。  
本手法は実行前に制約を強制する点が異なる。

\subsection{「それプロンプトでよくない？」}

プロンプトはソフト制約であり、必ず守られる保証がない。  
本手法ではハード制約として検証・補正を行う。

\subsection{「遅くならない？」}

一部オーバーヘッドはあるが、安全性とのトレードオフと考えている。  
また、軽量化や部分検証により改善可能である。
\end{document}